\subsection{Método de estimación, secuenciamiento y control}
Synaptix utiliza estimación ascendente por resultado: el diccionario fija contenido, exclusiones, recepción, responsable último, HH por perfil, recursos y supuestos; sus actividades asignan ejecutor, dedicación, lugar, calendario y predecesoras. La descomposición se detiene en resultados estimables, asignables y verificables. Las familias repetitivas comparten su ficha, mientras cada período o lote conserva código, responsable, esfuerzo y programación; el rango observado de 8--80 HH no es una obligación normativa. Un código identifica pertenencia a la EDT y no una fecha: se conservan los códigos históricos aunque cambie la ventana de ejecución. La etapa y el mes son atributos separados.

La ingeniería separa trabajo técnico y consolidación de Calidad; una aportación grande puede repartirse en subtareas entre personas competentes, sin sumar otra vez las HH. Lógica, persistencia e interfaz tienen integración asignada en la ficha. La consolidación de Calidad recibe los aportes previos; la validación final de marcha blanca permite en paralelo la verificación de datos por Calidad y la lista contractual por Proyecto, porque son resultados independientes. Las actividades y contribuciones individuales se identifican en \path{actividades_integradas.csv}.

La programación pasa por cinco controles: cobertura requisito--entregable--paquete; precedencias completas sin ciclos y puertas con responsable; CPM de avance/retroceso con holgura total y libre; escenario PERT de tres duraciones; y nivelación con calendario y capacidad. No se suma un porcentaje general de gestión o pruebas a trabajos ya estimados. Los servicios por turnos, las reservas y las esperas se muestran por separado. PERT no equivale a una probabilidad de entrega mientras no se sustenten distribuciones e independencia.

El calendario usa un escenario de inicio el 1 de diciembre de 2026 y feriados chilenos. La fecha contractual, los decretos de feriados, los catorce eventos náuticos y las campañas efectivas se incorporarán al control de cambios antes de aprobar la línea base. El método de este capítulo constituye la especificación que deberá recoger el capítulo metodológico de la oferta; no se atribuye aprobación o concordancia documental aún no realizada.


El esfuerzo de los núcleos se selecciona mediante UCP y se distribuye entre paquetes identificados, conservando una sola imputación del trabajo original. Las asignaciones y curvas corresponden a este alcance conciliado. La fecha de inicio sigue siendo una hipótesis explícita de calendario; la aprobación contractual de la línea base y la revisión humana técnica son controles futuros separados de estos cálculos.
